%% ma: L12-B04-0001
%% lop: 12
%% chuong: 1
%% bai: 4
%% dang: 12.4.1
%% loai: TN4
%% muc: NB
%% dap_an: "B"
%% nguon: "(NBV)-ĐỀ SỐ 3-MỖI NGÀY 1 ĐỀ THI 2025.pdf (D03, MỖI NGÀY 1 ĐỀ - Đề số 3), Phần I Câu 7"
%% kien_thuc: "Bài 4 (tâm đối xứng của đồ thị hàm phân thức bậc nhất trên bậc nhất, tiệm cận)"
%% trang_thai: cho-duyet
%% ghi_chu: "Đáp án gốc B đúng"
%% ly_do: "Dạng toán mới đề xuất 12.4.1 chờ giáo viên duyệt"
\begin{ex}
Đồ thị hàm số $y=f(x)=\dfrac{x-3}{x+2}$ có toạ độ của tâm đối xứng là
\choice{$(3;-2)$}{\True $(-2;1)$}{$\left(-2;-\dfrac32\right)$}{$\left(-\dfrac32;-2\right)$}
\loigiai{Đồ thị có tiệm cận đứng $x=-2$ và tiệm cận ngang $y=1$. Giao điểm hai tiệm cận $(-2;1)$ là tâm đối xứng của đồ thị.}
\end{ex}
